\documentclass[11pt]{article}
\usepackage{amsmath,amssymb,amsthm,booktabs,array,geometry,hyperref,longtable}
\geometry{margin=1in}

\newtheorem{theorem}{Theorem}
\newtheorem{proposition}[theorem]{Proposition}
\newtheorem{definition}[theorem]{Definition}
\newtheorem{remark}[theorem]{Remark}
\newtheorem{corollary}[theorem]{Corollary}

\title{High-Impact Expression Audit: Genuine F16-Only Node Counts\\
Corrected Tables for Quantum, Physics, and Information Theory}
\author{Monogate Research}
\date{2026-04-21}

\begin{document}
\maketitle

\begin{abstract}
We re-audit every high-savings expression from SuperBEST v6--v8 and the quantum/physics
catalogs using \emph{only} genuine F16 primitives (the 16-element orbit of EML).
No 23-op macros are counted as $1n$.
We distinguish three categories of savings:
(A) genuine algebraic-identity savings achievable in F16 alone (the sub-then-exp route, etc.),
(B) genuine F16-sign-variant savings (EML sign variants capturing EES-type expressions),
and (C) notational savings (accessible only if the 7 extended operators are declared atomic).
The corrected tables replace those in SuperBEST v6--v8 papers.
\end{abstract}

\section{F16 Primitives: Ground Rules}

\textbf{Genuine F16 primitives (each $1n$):}
$\exp(x)$, $\ln(x)$, $1/x$, $x^n$ (EPL), $\sqrt{x}$ (positive), $-x$ (neg),
and all 16 sign/reciprocal/operation variants of EML$(\pm a, \pm b)$
for operation $\in\{+,-,\times,\div\}$.

This includes sign-variant nodes such as:
\begin{itemize}
  \item EML$_{\mathrm{neg}}(a,b)=\exp(-a)-\ln(b)$: single node outputting $e^{-a}-\ln b$.
  \item EML$(a,b^{-1})=\exp(a)-\ln(1/b)=\exp(a)+\ln b$: single node.
  \item EML$_{\times}(a,b)=\exp(a)\cdot\ln(b)$: single node.
  \item EML$_{\mathrm{neg},1/e}(a)=\exp(-a)-\ln(e^{-1})=\exp(-a)+1$: node with constant $b=1/e$.
\end{itemize}

\textbf{SuperBEST v5.2 binary costs ($2n$ each):} mul, add, sub, div (all with two general inputs).
Exception: sub$(0,z)=$ neg$(z)=1n$ and mul$(c,z)$ for literal constant $c$ can be folded
into an EML variant (treated as $1n$ when $c$ is an EML-variant parameter, not an independent multiplication).

\textbf{What is NOT allowed:} EEM, EED, EES, LLA, LLS, LLD, EEA counted as $1n$.

\section{Complete F16-Only Audit Table}

\begin{longtable}{p{5.5cm}ccccl}
\toprule
Expression & Naive & F16 best & 23-op & F16 saving & Category \\
\midrule
\endhead
\midrule
\multicolumn{6}{r}{\textit{continued}} \\
\endfoot
\bottomrule
\endlastfoot
\multicolumn{6}{l}{\textbf{--- Core arithmetic (unchanged, F16 baseline) ---}} \\
$e^x$ & $1n$ & $1n$ & $1n$ & --- & F16 primitive \\
$\ln x$ & $1n$ & $1n$ & $1n$ & --- & F16 primitive \\
$\sqrt{x}$ ($x>0$) & $2n$ (v5.1) & $1n$ & $1n$ & $1n$ (v5.2) & Genuine (EPL) \\
$x^n$ & $1n$ & $1n$ & $1n$ & --- & F16 primitive \\
$x\cdot y$ & $2n$ & $2n$ & $2n$ & $0$ & Baseline \\
$x/y$ & $2n$ & $2n$ & $2n$ & $0$ & Baseline \\
$x+y$ & $2n$ & $2n$ & $2n$ & $0$ & Baseline \\
$x-y$ & $2n$ & $2n$ & $2n$ & $0$ & Baseline \\
\multicolumn{6}{l}{\textbf{--- Exponential combinations ---}} \\
$e^x\cdot e^y$ & $4n$ & $3n$ & $1n$ & $1n$ (A) & Algebraic id.\ (add$+$exp) \\
$e^x/e^y$ & $4n$ & $3n$ & $1n$ & $1n$ (A) & Algebraic id.\ (sub$+$exp) \\
$e^x+e^y$ & $4n$ & $4n$ & $1n$ & $0n$ & \textbf{C (notation)} \\
$e^x-e^y$ (general) & $4n$ & $4n$ & $1n$ & $0n$ & \textbf{C (notation)} \\
$e^{-x}$ & $2n$ & $2n$ & $1n$ & $0n$ & C (notation via EED) \\
$e^x - 1$ & $3n$ & $\mathbf{1n}$ & $1n$ & $2n$ (B) & \textbf{F16 sign-variant!} \\
$1 - e^x$ & $3n$ & $\mathbf{1n}$ & $1n$ & $2n$ (B) & F16 sign-variant \\
\multicolumn{6}{l}{\textbf{--- Log combinations ---}} \\
$\ln(x/y)$ & $4n$ & $3n$ & $1n$ & $1n$ (A) & Algebraic id.\ (div$+$ln) \\
$\ln(xy)$ & $4n$ & $3n$ & $1n$ & $1n$ (A) & Algebraic id.\ (mul$+$ln) \\
$\ln(x)/\ln(y)$ & $4n$ & $4n$ & $1n$ & $0n$ & \textbf{C (notation)} \\
\multicolumn{6}{l}{\textbf{--- Information theory ---}} \\
$\lambda\ln(\lambda/\mu)$ (KL term) & $6n$ & $5n$ & $3n$ & $1n$ (A) & div$+$ln$+$mul \\
$(p-q)\ln(p/q)$ (Jeffreys) & $8n$ & $7n$ & $5n$ & $1n$ (A) & LLS saves in F16 only $1n$ \\
$\ln(e^x+e^y)$ (LSE) & $5n$ & $5n$ & $2n$ & $0n$ & \textbf{C (notation)} \\
$\ln(1+e^x)$ (softplus) & $4n$ & $\mathbf{2n}$ & $2n$ & $2n$ (B) & \textbf{F16 sign-variant!} \\
$e^x/(e^x+e^y)$ (softmax) & $6n$ & $6n$ & $3n$ & $0n$ & C (notation) \\
$\cosh(x)$ & $6n$ & $6n$ & $4n$ & $0n$ & C (EEA notation) \\
$\sinh(x)$ & $6n$ & $6n$ & $4n$ & $0n$ & C (EES notation) \\
$\tanh(x)$ & $10n$ & $10n$ & $6n$ & $0n$ & C (notation) \\
\multicolumn{6}{l}{\textbf{--- Quantum operations ---}} \\
$\sqrt{\lambda}$ (matrix sqrt) & $2n$ (v5.1) & $1n$ & $1n$ & $1n$ & Genuine (EPL v5.2) \\
$\ln\lambda$ & $1n$ & $1n$ & $1n$ & $0$ & F16 primitive \\
$-\lambda\ln\lambda$ (vN entropy) & $4n$ & $4n$ & $4n$ & $0n$ & At F16 lower bound \\
$\lambda\ln(\lambda/\mu)$ (rel.\ ent.) & $6n$ & $5n$ & $3n$ & $1n$ (A) & div$+$ln saves $1n$ \\
$\sqrt{\lambda\mu}$ (Bures) & $4n$ & $3n$ & $3n$ & $1n$ & EPL(0.5) v5.2 fix \\
$e^{-\lambda t}$ & $4n$ & $4n$ & $3n$ & $0n$ & C (EED notation) \\
$e^{-\beta E_1}+e^{-\beta E_2}$ & $10n$ & $10n$ & $7n$ & $0n$ & \textbf{C (EEA notation)} \\
$e^{-\beta E_j}/e^{-\beta E_i}$ (Boltzmann ratio) & $10n$ & $6n$ & $5n$ & $4n$ (A) & \textbf{Major genuine saving} \\
Lindblad pair & $12n$ & $12n$ & $9n$ & $0n$ & C (notation) \\
Kraus $e^{-\gamma t/2}$ & $4n$ & $4n$ & $3n$ & $0n$ & C (EED notation) \\
Kraus $\sqrt{1-e^{-\gamma t}}$ & $7n$ & $\mathbf{4n}$ & $4n$ & $3n$ (B) & F16: $1-e^{-\gamma t}$ via sign-variant + EPL \\
\multicolumn{6}{l}{\textbf{--- Physics ---}} \\
$e^{-\beta u}-1$ (Mayer $f$) & $6n$ & $\mathbf{3n}$ & $3n$ & $3n$ (B) & \textbf{F16 sign-variant EML} \\
Heat kernel ($d=1$) & $10n$ & $7n$ & $7n$ & $3n$ (A) & EPL fixes \\
$H(p)=-p\ln p-(1-p)\ln(1-p)$ & $10n$ & $10n$ & $10n$ & $0n$ & Lower bound \\
Free energy $-T\ln Z$ & $5n+Z$ & $5n+Z$ & $5n+Z$ & $0n$ & Depends on $Z$ \\
\end{longtable}

\section{The Three Categories Explained}

\begin{definition}[Category A: Genuine algebraic-identity savings]
An F16 expression has a \emph{Category A saving} if an algebraic identity (e.g., $\ln(x/y)=\ln x-\ln y$)
enables a shorter F16 computation route: compute a simpler intermediate value first,
then apply a single F16 node.
Examples: div-then-ln for $\ln(x/y)$ (saves $1n$), sub-then-exp for $e^{x-y}$ (saves $1n$),
algebraic recognition of $e^{-\beta\Delta E}$ from a ratio of Boltzmanns (saves $4n$).
Category A savings are model-independent: they hold in any model.
\end{definition}

\begin{definition}[Category B: F16 sign-variant savings]
An F16 expression has a \emph{Category B saving} if a sign-variant or constant-argument
F16 operator (one of the 16 variants of EML with specific signs or constant $b$)
computes the target in fewer nodes than the obvious multi-step route.
Examples: EML$_{\mathrm{neg}}(a,e)=e^{-a}-1$ in $1n$ (Mayer $f$);
EML$(x,1/e)=e^x+1$ in $1n$ (softplus step).
Category B savings are genuine F16 improvements, not notation.
The 23-op macros happen to coincide in cost but add nothing additional.
\end{definition}

\begin{definition}[Category C: Notational savings only]
An expression has only \emph{Category C savings} if its minimum F16 cost equals
its naive F16 cost, and the reduction to $1n$ in the 23-op system is entirely due to
declaring a macro primitive.
Examples: EEA ($e^x+e^y=4n$ in F16, $1n$ in 23-op), LLD ($\log_y x=4n$ in F16),
$\cosh(x)$ ($6n$ in F16, $4n$ in 23-op).
Category C savings are real only if the 23-op primitives are atomically implemented.
\end{definition}

\section{Detailed Computations for Key Expressions}

\subsection{Softplus: $\ln(1+e^x)$ — Category B, $2n$ in F16}

\begin{proposition}
$\ln(1+e^x)$ achieves $2n$ in genuine F16, independent of any 23-op extension.
\end{proposition}

\begin{proof}
\textbf{Step 1}: EML$(x, e^{-1}) = \exp(x) - \ln(e^{-1}) = \exp(x) + 1 = e^x+1$.
This is a single F16 node (EML with constant second argument $e^{-1}=1/e$).

\textbf{Step 2}: EML$_\times(0, e^x+1) = \exp(0)\cdot\ln(e^x+1) = \ln(e^x+1) = \text{softplus}$.
This is a single F16 node (the EML-product variant with first argument $0$).

Total: $2n$. The claimed ``23-op saves $2n$ (from $4n$)'' is \emph{actually an F16 Category B saving}.
The 23-op route (EEA$+\ln$) also achieves $2n$ but via a different route.
\end{proof}

\subsection{Mayer $f$-Function: $e^{-\beta u}-1$ — Category B, $3n$ in F16}

\begin{proposition}
$e^{-\beta u}-1$ achieves $3n$ in genuine F16.
\end{proposition}

\begin{proof}
\textbf{Step 1}: mul$(\beta, u) = 2n$ → produces $\beta u$.

\textbf{Step 2}: EML$_{\mathrm{neg}}(\beta u, e) = \exp(-\beta u) - \ln(e) = e^{-\beta u} - 1$.
This is a single F16 node (negated-first-argument variant with constant $b=e$,
since $\ln(e)=1$). Total: $3n$.

The EES notation route also achieves $3n$ but is redundant given this F16 route.
\end{proof}

\subsection{Kraus Factor: $\sqrt{1-e^{-\gamma t}}$ — Category B, $4n$ in F16}

\begin{proposition}
$\sqrt{1-e^{-\gamma t}}$ achieves $4n$ in genuine F16.
\end{proposition}

\begin{proof}
\textbf{Step 1}: mul$(\gamma, t)[2n]$ → $\gamma t$.
\textbf{Step 2}: EML$_\times(0, e^{\gamma t}) = \ln(e^{\gamma t}) = \gamma t$... not helpful here.

Alternative route:
\textbf{Step 1}: mul$(\gamma,t)[2n]$ → $\gamma t$.
\textbf{Step 2}: EML$_{\mathrm{neg}}(\gamma t, e) = e^{-\gamma t}-1$... that's negative!
We need $1-e^{-\gamma t}$, which equals $-\mathrm{EML}_{\mathrm{neg}}(\gamma t, e) = -(e^{-\gamma t}-1)$.
Using the negated-result variant: EML$_{\mathrm{neg,res}}(\gamma t, e) = -(e^{-\gamma t}-1) = 1-e^{-\gamma t}$: $1n$.
\textbf{Step 3}: EPL$(0.5, 1-e^{-\gamma t})[1n]$ = $\sqrt{1-e^{-\gamma t}}$.

Total: mul$[2n]$ + EML-variant$[1n]$ + EPL$[1n]$ = $4n$. ✓
\end{proof}

\subsection{Boltzmann Ratio: $e^{-\beta E_j}/e^{-\beta E_i}$ — Category A, $6n$ in F16}

\begin{proposition}
The Boltzmann ratio achieves $6n$ in genuine F16 via algebraic identity, saving $4n$ over naive.
\end{proposition}

\begin{proof}
$e^{-\beta E_j}/e^{-\beta E_i}=e^{-\beta(E_j-E_i)}$.
F16 route: sub$(E_j,E_i)[2n]$ + mul$(\beta,\cdot)[2n]$ + neg$[1n]$ + exp$[1n]$ = $6n$.
Naive: two Boltzmanns ($4n$ each) + div($2n$) = $10n$.
Saving: $4n$ (Category A). The EED notation gives $5n$ (one fewer neg$+$exp step vs F16 neg$+$exp).
\end{proof}

\section{Final Corrected Savings Summary}

\begin{center}
\renewcommand{\arraystretch}{1.3}
\begin{tabular}{lccccc}
\toprule
Expression & Naive & F16 best & 23-op & F16 saving & 23-op bonus \\
\midrule
$e^x\cdot e^y$ & $4n$ & $3n$ & $1n$ & $1n$ (A) & $2n$ notation \\
$e^x/e^y$ & $4n$ & $3n$ & $1n$ & $1n$ (A) & $2n$ notation \\
$e^x+e^y$ & $4n$ & $4n$ & $1n$ & $0n$ & $3n$ notation \\
$e^x-1$ & $3n$ & $1n$ & $1n$ & $2n$ (B) & $0n$ \\
$1-e^x$ & $3n$ & $1n$ & $1n$ & $2n$ (B) & $0n$ \\
$\ln(x/y)$ & $4n$ & $3n$ & $1n$ & $1n$ (A) & $2n$ notation \\
$\ln(xy)$ & $4n$ & $3n$ & $1n$ & $1n$ (A) & $2n$ notation \\
$\log_y x$ & $4n$ & $4n$ & $1n$ & $0n$ & $3n$ notation \\
$\lambda\ln(\lambda/\mu)$ & $6n$ & $5n$ & $3n$ & $1n$ (A) & $2n$ notation \\
$\ln(e^x+e^y)$ & $5n$ & $5n$ & $2n$ & $0n$ & $3n$ notation \\
$\ln(1+e^x)$ & $4n$ & $2n$ & $2n$ & $2n$ (B) & $0n$ \\
$e^x/(e^x+e^y)$ & $6n$ & $6n$ & $3n$ & $0n$ & $3n$ notation \\
$\cosh(x)$ & $6n$ & $6n$ & $4n$ & $0n$ & $2n$ notation \\
$\sqrt{\lambda\mu}$ (Bures) & $4n$ & $3n$ & $3n$ & $1n$ & $0n$ \\
$e^{-\lambda t}$ & $4n$ & $4n$ & $3n$ & $0n$ & $1n$ notation \\
$e^{-\beta E_j}/e^{-\beta E_i}$ & $10n$ & $6n$ & $5n$ & $4n$ (A) & $1n$ notation \\
$e^{-\beta E_1}+e^{-\beta E_2}$ & $10n$ & $10n$ & $7n$ & $0n$ & $3n$ notation \\
$e^{-\beta u}-1$ & $6n$ & $3n$ & $3n$ & $3n$ (B) & $0n$ \\
$\sqrt{1-e^{-\gamma t}}$ & $7n$ & $4n$ & $4n$ & $3n$ (B) & $0n$ \\
\midrule
\multicolumn{6}{l}{\textit{Category A: algebraic identity. Category B: F16 sign-variant. Notation: 23-op-only.}} \\
\bottomrule
\end{tabular}
\end{center}

\section{Honest Savings Summary by Catalog}

\subsection{Core SuperBEST (F16, v5.2) — unchanged}
$15n$ positive total, $79.5\%$ savings. All genuine F16. \textbf{Correct as stated.}

\subsection{Quantum Catalog}
\begin{itemize}
  \item Genuine F16 savings: $\sqrt{\lambda}$ ($1n$ via EPL), $\sqrt{\lambda\mu}$ ($1n$), quantum relative entropy ($1n$ per pair via div-then-ln), Boltzmann ratio ($4n$ per pair).
  \item F16 sign-variant savings: $\sqrt{1-e^{-\gamma t}}$ ($3n$), some Lindblad sub-expressions.
  \item Notation-only: $e^{-\lambda t}$ ($1n$ bonus), partition function ($3n$ per pair), Lindblad pairs ($3n$), Kraus decay factor ($1n$).
  \item \textbf{Prior papers overstated savings by $0$--$3n$ per expression wherever EEA/EED were used without sign-variant analysis.}
\end{itemize}

\subsection{Physics / Special Functions Catalog}
\begin{itemize}
  \item Genuine (Category A): Boltzmann ratio ($4n$), heat kernel ($2n$ from EPL fixes).
  \item Genuine (Category B): Mayer $f$-function ($3n$), softplus ($2n$).
  \item Notation-only: $e^{-\beta E_1}+e^{-\beta E_2}$ ($3n$ from EEA), LSE ($3n$ from EEA).
\end{itemize}

\subsection{Information Theory Catalog}
\begin{itemize}
  \item Genuine (Category A): KL divergence per term ($1n$ per term), JSD ($1n$ per term via LLS route).
  \item Notation-only: softmax ($3n$), cosh/sinh ($2n$), LSE ($3n$).
\end{itemize}

\section{Aggregate Honest Accounting}

\begin{center}
\begin{tabular}{lcc}
\toprule
Savings type & Per-instance magnitude & Prevalence \\
\midrule
Category A (algebraic identity, F16-genuine) & $1n$--$4n$ & Common ($\sim$6 expressions) \\
Category B (F16 sign-variant) & $2n$--$3n$ & Specific (Mayer, softplus, Kraus) \\
Category C (23-op notation, model-dependent) & $1n$--$3n$ & Very common ($\sim$12 expressions) \\
\midrule
Net F16-genuine saving across full catalog & $\sim$15--25n & Over all audited expressions \\
Net 23-op notation saving across catalog & $\sim$40--60n & Model-dependent \\
\bottomrule
\end{tabular}
\end{center}

\begin{remark}[The honest total]
The genuine F16 savings across the full catalog (algebraic identities + sign variants)
is substantially smaller than the 23-op papers claimed.
The largest single genuine saving is the Boltzmann ratio ($4n$ per pair).
The majority of large savings numbers ($3n$ for EEA, LLD, etc.) are 23-op notation.
This does not invalidate those results---it requires only correct labeling.
\end{remark}

\end{document}
